\chapter{Security and privacy}
\label{ch:security}

Yap handles three kinds of sensitive data: where people are, what they look like, and whom they talk
to. This chapter describes how the system protects each, starting from what it assumes about the client.
\Cref{fig:security} shows every path from a phone to the data and the rule that guards it.

\begin{figure}[ht]
  \centering
  % Security boundaries: every path from the phone to the data and the rule on it (Figure: security).
\begin{tikzpicture}[
    svc/.style={box, serverbox, font=\sffamily\scriptsize, text width=24mm, minimum height=10mm, align=center},
    dbbox/.style={box, font=\sffamily\scriptsize, align=center},
    gate/.style={font=\sffamily\tiny, text=accentdark, inner sep=1.2pt, align=center},
  ]
  \node[box, clientbox, font=\sffamily\scriptsize, text width=24mm, minimum height=62mm, align=center] (phone) at (1.25,-2.5)
    {\textbf{\footnotesize Phone}\\[4pt]the app, or any client built with the public key\\[4pt]session token: role \texttt{authenticated}\\[4pt]{\color{alert}no token: role \texttt{anon}, which has no privileges on the app's tables or functions}};

  \node[svc] (rpc) at (5.1,0.5) {Data API: \texttt{rpc}};
  \node[svc] (rest) at (5.1,-1.6) {Data API: tables};
  \node[svc] (storage) at (5.1,-3.4) {Storage API};
  \node[svc] (rt) at (5.1,-5.2) {Realtime};

  \begin{scope}[on background layer]
    \node[layer, fit={(7.3,1.7) (15.75,-6.0)}, inner sep=0pt] (db) {};
  \end{scope}
  \node[layerlabel, anchor=north east] at (db.north east) {POSTGRES};

  \node[dbbox, accentbox, text width=30mm] (fns) at (9.65,0.5) {\textbf{\texttt{public} functions}\\\texttt{SECURITY DEFINER}\\\texttt{search\_path = ''}\\check, then write as owner\\{\color{muted}helpers in \texttt{private}, not exposed}};
  \node[dbbox, text width=28mm, minimum height=27mm] (tables) at (13.95,-0.8) {\textbf{\texttt{public} tables}\\[3pt]row-level security on every table, deny by default\\[3pt]column grants on \texttt{profiles}};
  \node[dbbox, text width=24mm] (objects) at (13.7,-3.4) {\textbf{\texttt{storage.objects}}\\bucket \texttt{photos}, private};

  \foreach \s in {rpc, rest, storage} {\draw[flow] (phone.east |- \s) -- (\s);}
  \draw[flow] (rt) -- (phone.east |- rt);

  \draw[flow] (rpc) -- node[gate, above=2pt] {\texttt{EXECUTE}:\\signed in} (fns);
  \draw[flow] (fns.east |- 0,0.2) -- node[gate, above=2pt] {as owner} (tables.west |- 0,0.2);
  \draw[flow] (rest) -- (tables.west |- rest);
  \node[gate, anchor=south west, text width=40mm, align=flush left] at (6.75,-1.55) {policies: own profile; requests and messages the user takes part in; own blocks (select, delete)};
  \draw[flow] (storage) -- (objects);
  \node[gate, anchor=south west, text width=44mm, align=flush left] at (6.75,-3.35) {write: own folder only; read: own folder, or \texttt{private.can\_see\_photo}};
  \draw[flow] (15.45,0 |- tables.south) -- (15.45,-5.2) -- node[gate, above=2pt] {changes to requests and messages, only rows the subscriber may select} (rt.east);
\end{tikzpicture}

  \caption{Security boundaries. Writes go through functions that check before they act; direct table and
    storage access is limited by policies; Realtime delivers only rows the subscriber could select.}
  \label{fig:security}
\end{figure}

\section{Trust model}

The server trusts nothing the client says about identity or state. The app's address and publishable
key are in the app bundle and must be assumed public, so anyone can build a client that calls the same
API with a valid user session. Every rule therefore lives in the database, where such a client cannot
bypass it:

\begin{itemize}
  \item Identity comes only from the session token, through \code{auth.uid()}.
  \item State, such as whether a user is live, chatting or waiting on a request, is read from the
    database, never taken from arguments.
  \item Proximity is computed from stored positions, never from a distance the client supplies.
  \item Every write passes through a database function, except three the policies confine to the user's
    own data: the profile's client-writable columns, deleting one's own \code{blocks} rows (unblocking),
    and photo files in one's own Storage folder. Each function checks what its action needs; blocks and
    rate limits apply where \cref{tab:functions} lists them, so \code{update_note}, \code{leave_event}
    and \code{end_chat}, for example, check neither.
\end{itemize}

The server must accept two inputs it cannot verify: the position the phone reports, and the GPS accuracy,
which widens the nearby radius by at most 50 metres. The relocation limits bound how a live position may
change (\cref{sec:relocation-limits}); the position sent to join an event is checked only against the
event's distance.

\section{Authentication and sessions}

Accounts are Supabase Auth users with an email address and a password of at least eight characters.
Every new address is confirmed with an emailed 8-digit code before the account can sign in, and a
forgotten password is reset with another emailed code. Supabase limits how often codes are sent; the app
reports those limits as sentences, except when creating an account, which looks the same whether or not
the email already has an account (\cref{sec:flow-account}).

The Supabase client keeps the session in AsyncStorage, refreshes its access token automatically, and does
not look for sessions in URLs. AsyncStorage is used rather than the platform's secure store because the
session JSON is larger than the roughly 2\,KB a secure-store entry holds. Before joining any Realtime
channel the client passes the current token to Realtime, so every channel joins as the signed-in user and
receives only the rows that user may read. On sign-out the app closes every channel, forgets every
signed URL and clears the photos cached on the phone, so the next account to sign in on the same phone
sees nothing of the last one.

\section{Row-level security and grants}
\label{sec:rls}

Row-level security is enabled on every application table and on the three private tables,
\code{private.location_changes}, \code{private.event_posts} and \code{private.blocked_words}. The
anonymous role has no access to any of the app's tables and cannot call any of its functions. \Cref{tab:rls} lists what a
signed-in user may do directly; everything else happens inside a database function.

\begin{table}[ht]
  \centering
  \caption{Direct access for signed-in users. A dash means no access.}
  \label{tab:rls}
  \tablefont
  \begin{tabularx}{\textwidth}{@{}L{28mm}L{34mm}L{30mm}Y@{}}
    \toprule
    \textsf{\textbf{Table}} & \textsf{\textbf{Select}} & \textsf{\textbf{Insert, update}} & \textsf{\textbf{Delete}} \\
    \midrule
    \code{profiles} & own row & own row; editable columns only; \code{is_adult} must be true on insert & --- \\
    \code{availability} & --- & --- & --- \\
    \code{chat_requests} & rows the user sent or received & --- & --- \\
    \code{messages} & chats the user is in and events the user has joined, minus senders the user blocked & --- & --- \\
    \code{events}, \code{event_members} & --- & --- & --- \\
    \code{blocks} & blocks the user made & --- & blocks the user made \\
    \code{reports} & --- & --- & --- \\
    \code{location_changes}, \code{event_posts}, \code{blocked_words} & --- & --- & --- \\
    \code{storage.objects} (\code{photos}) & own folder, or a photo of someone the user can see & directly in own folder, at most 24 files & own folder \\
    \bottomrule
  \end{tabularx}
\end{table}

Column grants close a gap that row-level security alone would leave: a user may update their own profile
row, but only its editable columns. \code{verified} and \code{status} cannot be written by any client, and
\code{id} and \code{is_adult} cannot be changed after the row is created.

Every function is hardened in the same way. \code{SECURITY DEFINER} functions set an empty
\code{search_path}, so a caller cannot make them resolve a name to an object of the caller's choosing.
Execution is revoked from \code{public} and \code{anon}. Helpers live in a schema the Data API does not
expose. A platform helper that Supabase installs, \code{public.rls_auto_enable}, is also revoked from
every API role, because a \code{SECURITY DEFINER} function should not be callable by them even when it
can do nothing outside its intended trigger.

\section{What the functions enforce}

\Cref{tab:functions} lists every check. They fall into seven groups:

\begin{description}
  \item[Identity.] Every function acts as \code{auth.uid()}; a user cannot request, block or report
    themselves.
  \item[Profile standing.] Going live, creating and joining events require an active adult profile;
    only active users can be found or sent a request.
  \item[State.] A request can be answered only while it is \code{pending} and addressed to the caller; a
    chat can be ended or written to only while it is \code{accepted} and the caller takes part in it;
    one request may wait per sender and one chat may run per user; only its host can delete an event.
  \item[Proximity.] Nearby lists and requests require both users to be live within 100 to 150 metres,
    and so does seeing a stranger's photos; joining an event requires being within 5 kilometres of it.
  \item[Blocks.] Any block between two users, in either direction, stops every interaction between
    them (\cref{sec:blocking-reporting}).
  \item[Limits.] Lengths, durations, start times and rates (\cref{tab:rate-limits}).
  \item[Content.] Typed text must show something and carry no control characters, direction overrides or
    piles of accents (check constraints), and text other people see may not contain a word from the word
    filter's list (\cref{sec:word-filter}).
\end{description}

\section{Location privacy}
\label{sec:location-privacy}

\subsection{How a position is handled}

A live position follows a short, fixed path:

\begin{enumerate}
  \item The app asks for \emph{while using the app} permission only, never background location. It reads a
    position when the user taps \ui{Go live} or \ui{Join}, and when the \ui{Map} tab centres on the user. It
    does not follow the user's movements.
  \item \code{go_green} stores the position in the user's \code{availability} row, which no client can read.
  \item Database functions use the position only to compare distances. No function returns it, and no
    function returns a distance: \code{nearby} returns person cards, and the cards contain no location.
  \item \code{go_offline} deletes the row the moment the user goes offline, and the minute job deletes it
    when the session ends. No copy is kept. \code{private.location_changes} keeps the time of each
    location change for ten minutes, but never the place.
\end{enumerate}

The positions sent to \code{events_near} (the centre of the map, usually the user's position) and to
\code{join_event} are used for one query each and are not stored. Event locations
are the exception to all of this by design: they are places hosts choose to publish, and the map shows
them to anyone within the chosen radius.

The nearby radius is deliberately coarse. A user appears to others only as ``nearby'', within 100 metres
plus the viewer's reported accuracy and never beyond 150 metres, and moves of 150 metres or less do not
change the stored position at all (\cref{sec:relocation-limits}), so the list cannot be used to track
someone's small movements.

\subsection{Relocation limits}
\label{sec:relocation-limits}

A modified client could call \code{go_green} with arbitrary coordinates and then \code{nearby}, and so
learn who is live at places the user is not. \code{go_green} limits this probing in three ways:

\begin{itemize}
  \item \textbf{GPS wobble is ignored.} If the user is live and the new position is within 150 metres of
    the stored one, the stored position and its time stay as they were. Small steps therefore cannot add
    up to a large move.
  \item \textbf{Speed.} While the user is live, a larger move must not exceed 150 metres plus walking pace,
    10 km/h, over the time since the stored position was set (\cref{lst:speed}).
  \item \textbf{Frequency.} Starting a session, or moving more than 150 metres, counts as a location
    change. At most three are allowed in any ten minutes.
\end{itemize}

\begin{lstlisting}[language=pgsql, caption={The speed check in \code{go_green}. \code{moved_m} is the distance from the stored position.}, label={lst:speed}]
-- 10 km/h is 10000 m per 3600 s, counted from when your live position was set
allowed_m := 150 + 10000.0 / 3600 * extract(epoch from now() - previous.started_at);
if moved_m > allowed_m then raise exception 'moving too fast'; end if;
\end{lstlisting}

\subsection{The delete-on-offline trade-off}
\label{sec:trade-off}

The go-live screen promises that the location is deleted when the user stops. Keeping that promise means
the server holds no position once a user is offline, and so the speed check can apply only within one
live session. A modified client can therefore go offline, go live somewhere else, and repeat, limited only
by the frequency rule: three location changes in ten minutes, or about eighteen an hour. The alternative,
keeping each user's last position after they go offline so that the speed rule could span sessions, would
have broken the promise. The project chose to keep the promise and accept this residual risk; no cooldown
was added.

\section{Photo privacy}
\label{sec:photo-privacy}

Photos are the most identifying data in Yap, and their protection has four layers.

\paragraph{A private bucket} The \code{photos} bucket is private, so no photo has a public address.
Clients see a photo only through a signed URL, valid for one hour.

\paragraph{Own folder only} Storage policies let a user upload, overwrite and delete only inside the
folder named after their own user ID. Check constraints on \code{profiles} require the profile picture and
every further photo to be a file directly inside the owner's folder, with no deeper path and no
\code{..}, and allow at most six further photos. A profile therefore cannot show someone else's face.

\paragraph{Visible people only} Reading an object, which signing a URL and listing a folder both require,
is allowed for the user's own folder and otherwise only when \code{private.can_see_photo} agrees. It
allows a photo only if all of the following hold:

\begin{itemize}
  \item the file is on the owner's current profile (the profile picture, its 320-pixel copy, or one of
    the further photos), never an older file left in the folder;
  \item neither user has blocked the other;
  \item and one of these is true: the owner is on the user's nearby list right now (owner active, live and
    open, and both near each other); a request between them is waiting, or they are chatting or met; or
    the file is the owner's profile picture and the owner hosts an event the map shows.
\end{itemize}

\paragraph{Short-lived URLs} The policy is checked when a URL is signed. A URL signed while its owner was
visible stays usable until it expires, at most one hour later, even if the owner stops being visible in
the meantime.

\section{What other users can see}

\Cref{tab:exposure} sets out what another user can learn about a person, and when.

\begin{table}[ht]
  \centering
  \caption{What other users can see.}
  \label{tab:exposure}
  \tablefont
  \begin{tabularx}{\textwidth}{@{}L{42mm}Y@{}}
    \toprule
    \textsf{\textbf{Data}} & \textsf{\textbf{Visible to others}} \\
    \midrule
    First name, last initial & On the person card, in event cards (host) and with every message sent \\
    Last name, email address & Never \\
    Gender & On the person card, only if the owner chose to show it and did not choose \ui{Prefer not to say} \\
    Introduction, interests & On the person card \\
    Photos & Through signed URLs, only while the owner is visible to the viewer (\cref{sec:photo-privacy}) \\
    Live note & On the person card while the owner is live; through \code{request_person} only while the request or chat is still going \\
    Live position, distance & Never \\
    Event location & To anyone whose map radius covers it, by design \\
    Being live & Only to other live users nearby, and only while the owner is \code{open} \\
    Blocks and reports & Never; a reported user is not told who reported them \\
    \bottomrule
  \end{tabularx}
\end{table}

The 18+ confirmation is the user's own declaration; the app does not verify age. \code{verified} is
reserved for the selfie check, which is switched off (\cref{sec:feature-flags}).

\section{Blocking and reporting}
\label{sec:blocking-reporting}

A block is stored as one directed row but enforced in both directions. The helper
\code{private.blocked_between} is checked by \code{nearby}, \code{send_request}, \code{request_person},
\code{send_chat_message}, \code{events_near}, \code{join_event} and \code{can_see_photo}, so once either
user blocks the other they no longer appear to each other, cannot send each other requests or chat
messages, cannot load each other's photos, and do not see each other's events. \code{block_user} also
cancels any waiting request and any chat between the two at once. Event group chats are shared spaces: a
blocked user can still write in an event both have joined, but the user who blocked them does not receive
those messages, because the \code{messages} select policy removes them, and the app filters them again.
Unblocking deletes the row, and the person can appear again at the next refresh.

A report records the reason and details. When three different users have reported the same profile, the
profile becomes \code{hidden}: it disappears from nearby lists, cannot receive requests, and its owner can
no longer go live or create or join events, until the owner reviews it. The owner reviews open reports at
least once every 24 hours, as the app stores require, from the Supabase SQL editor with the queries in
\code{docs/legal/moderation.md}, and records each outcome in \code{reviewed_at} and \code{action}. A
profile that broke the guidelines is suspended with \code{private.suspend_user}, which also ends
everything the person is part of (\cref{sec:moderation}); a hidden profile found harmless is set back to
\code{active}. There is no moderator screen in the app. The community guidelines, which every user agrees
to at sign-up, say what is not allowed and that reports are reviewed within 24 hours.

\section{Rate limits and abuse controls}

\Cref{tab:rate-limits} collects the limits that bound what one account can do.

\begin{table}[ht]
  \centering
  \caption{Rate limits.}
  \label{tab:rate-limits}
  \tablefont
  \begin{tabularx}{\textwidth}{@{}L{44mm}L{48mm}Y@{}}
    \toprule
    \textsf{\textbf{Action}} & \textsf{\textbf{Limit}} & \textsf{\textbf{Enforced by}} \\
    \midrule
    Sending requests & 5 per sender in any 10 minutes & \code{send_request} \\
    Waiting requests & 1 per sender at a time & \code{send_request}; partial unique index \\
    Sending messages & 30 per sender per minute, both chat kinds & \code{send_chat_message}, \code{send_event_message} \\
    Creating events & 5 per host in any 24 hours, deleted events included; exact for calls sent at once & \code{create_event} (a per-host lock); \code{private.event_posts} \\
    Reporting & 1 report per person per day; 10 people per day; exact for calls sent at once & \code{report_user} (a per-reporter lock) \\
    Storing photos & 24 files per user; best-effort, as uploads sent at once can go a little past it & storage insert policy \\
    Changing location & 3 changes in any 10 minutes & \code{go_green} \\
    Moving while live & 150\,m plus 10\,km/h since the position was set & \code{go_green} \\
    Email codes & 1 per address per minute on the hosted project; 1 per second on the local stack & Supabase Auth: a dashboard setting; \code{max_frequency} in \code{supabase/config.toml} \\
    Auth requests & Supabase's per-project limits & Supabase Auth \\
    \bottomrule
  \end{tabularx}
\end{table}

\section{Keys and secrets}

The app contains two configuration values, the project URL and the publishable key, both in
\code{src/config/supabase.ts}. Both are public by design: the key grants only what row-level security and
the database functions allow. No secret key, service-role key or SMTP credential is in the repository or
the app; they exist only in the Supabase dashboard. The Android build adds a third value, the Google Maps
key, which ships inside the Android app and is restricted in Google Cloud to the app's package and signing
certificate; it is kept out of the repository as an EAS environment variable.
